\documentclass{article}
\usepackage[paperwidth=36in,paperheight=24in,hmargin=1.4in,vmargin=1in]{geometry}

\usepackage{graphicx}
\usepackage{enumitem}
\pagenumbering{gobble}
\usepackage[utf8]{inputenc}
\usepackage{emoji}
\usepackage{tikz}
\usetikzlibrary{positioning,fit,calc}
\usepackage{tikz-qtree}
\usepackage{amsmath}
\usepackage{booktabs}
\usepackage{tabularx}
\usepackage{hyperref}
\hypersetup{
    colorlinks=true,
    linkcolor=blue,  % Change link color to blue
    urlcolor=blue    % Change URL color to blue
}
\usepackage{newpxtext,newpxmath}
\usepackage{eulervm}
\makeatletter
\def\txds@@scale{1}
\input{ot1tx-ds.fd}
\makeatother
\usepackage[bb=tx]{mathalpha}

\newcommand{\titlefont}{\fontsize{55}{55}\selectfont\setlength{\parskip}{1\baselineskip}}
\newcommand{\sectionfont}{\fontsize{34}{34}\selectfont\setlength{\parskip}{1\baselineskip}}
\newcommand{\nfont}{\fontsize{21}{22}\selectfont\setlength{\parskip}{1\baselineskip}}
\newcommand{\tikzfont}{\fontsize{11}{13}\selectfont\setlength{\parskip}{1.2\baselineskip}}
\newcommand{\agdafont}{\fontsize{18}{19}\selectfont\setlength{\parskip}{1\baselineskip}}
\newcommand{\jsonfont}{\fontsize{18}{19}\selectfont\setlength{\parskip}{1\baselineskip}}

\newcommand{\light}[1]{\textcolor{gray!90}{#1}}
\newcommand{\ie}{\textit{i.e.},}
\newcommand{\ia}{\textit{i.a.},}



%%%%%%%%%%
% Minted %
%%%%%%%%%%
\usepackage{minted}
\setminted[json]{fontsize=\jsonfont}
\setminted[agda]{fontsize=\agdafont}

%%%%%%%%
% Agda %
%%%%%%%%
\usepackage[conor]{agda}
\AtBeginEnvironment{code}{\agdafont}
\usepackage{newunicodechar}
\newunicodechar{ℕ}{\ensuremath{\mathbb{N}}}
\newunicodechar{λ}{\ensuremath{\mathnormal\lambda}}
\newunicodechar{→}{\ensuremath{\mathnormal\to}}
\newunicodechar{∀}{\ensuremath{\mathnormal\forall}}
\newunicodechar{≡}{\ensuremath{\mathnormal\equiv}}
\newunicodechar{⋯}{\ensuremath{\mathnormal\dots}}
\newcommand\AC[1]{\AgdaInductiveConstructor{\agdaFontSize#1}}
\newcommand\AD[1]{\AgdaDatatype{\agdaFontSize#1}}
\newcommand\AF[1]{\AgdaFunction{\agdaFontSize#1}}
\newcommand\AS[1]{\AgdaSymbol{\agdaFontSize#1}}
\newcommand\AB[1]{\AgdaBound{\agdaFontSize#1}}
\newcommand\AK[1]{\AgdaKeyword{\agdaFontSize#1}}

\newcommand{\fanode}{@}
\newcommand{\sanode}{$\to$}
\newcommand{\dpnode}{$\Pi$}



\begin{document}
	% Title
	\begin{minipage}{\textwidth}
		\begin{minipage}{0.5\textwidth}
			\titlefont
			Learning Structure-Aware Representations \\of Dependent Types
		\end{minipage}%
		\begin{minipage}{0.3\textwidth}
			\nfont
			Konstantinos Kogkalidis\\
			Orestis Melkonian\\
			Jean-Philippe Bernardy\\
		\end{minipage}\hfill
		\begin{minipage}{0.15\textwidth}
			\includegraphics[width=\textwidth]{NeurIPS_logo.png}
		\end{minipage}
	\end{minipage}
	\vspace{0.5in}

	\begin{minipage}[t]{0.29\textwidth}
	{\sectionfont\textbf{Background}}
	\nfont
	
	\textbf{Context}. Dependent type theories are formal languages used for defining mathematical objects and reasoning about their properties. Dependently-typed programming languages equate proofs with programs, facilitating theorem proving and formal verification. Here's a tiny program in \textbf{Agda}, proving that the addition of naturals is commutative:
	
	\hfill\begin{minipage}{0.9\textwidth}
		\begin{code}%
\>[0]\AgdaKeyword{open}\AgdaSpace{}%
\AgdaKeyword{import}\AgdaSpace{}%
\AgdaModule{Relation.Binary.PropositionalEquality}\AgdaSpace{}%
\AgdaKeyword{using}\AgdaSpace{}%
\AgdaSymbol{(}\AgdaOperator{\AgdaDatatype{\AgdaUnderscore{}≡\AgdaUnderscore{}}}\AgdaSymbol{;}\AgdaSpace{}%
\AgdaInductiveConstructor{refl}\AgdaSymbol{;}\AgdaSpace{}%
\AgdaFunction{cong}\AgdaSymbol{;}\AgdaSpace{}%
\AgdaFunction{trans}\AgdaSymbol{)}\<%
\end{code}
\begin{code}[hide]%
\>[0]\AgdaComment{--\ \AgdaUnderscore{}≡\AgdaUnderscore{}\ \ \ :\ ∀\ \{A\ :\ Set\}\ →\ A\ →\ A\ →\ Set\ a}\<%
\\
\>[0]\AgdaComment{--\ refl\ \ :\ ∀\ \{A\ :\ Set\}\ \{x\ :\ A\}\ →\ x\ ≡\ x}\<%
\\
\>[0]\AgdaComment{--\ cong\ \ :\ ∀\ \{A\ B\ :\ Set\}\ (f\ :\ A\ →\ B)\ \{x\ y\ :\ A\}\ →\ x\ ≡\ y\ →\ f\ x\ ≡\ f\ y}\<%
\\
\>[0]\AgdaComment{--\ trans\ :\ ∀\ \{A\ :\ Set\}\ \{i\ j\ k\ :\ A\}\ →\ i\ ≡\ j\ →\ j\ ≡\ k\ →\ i\ ≡\ k}\<%
\end{code}
\begin{minipage}{.3\linewidth}
\begin{code}%
\>[0]\AgdaKeyword{data}\AgdaSpace{}%
\AgdaDatatype{ℕ}\AgdaSpace{}%
\AgdaSymbol{:}\AgdaSpace{}%
\AgdaPrimitive{Set}\AgdaSpace{}%
\AgdaKeyword{where}\<%
\\
\>[0][@{}l@{\AgdaIndent{0}}]%
\>[2]\AgdaInductiveConstructor{zero}%
\>[8]\AgdaSymbol{:}\AgdaSpace{}%
\AgdaDatatype{ℕ}\<%
\\
%
\>[2]\AgdaInductiveConstructor{suc}%
\>[8]\AgdaSymbol{:}\AgdaSpace{}%
\AgdaDatatype{ℕ}\AgdaSpace{}%
\AgdaSymbol{→}\AgdaSpace{}%
\AgdaDatatype{ℕ}\<%
\end{code}
\end{minipage}

\begin{minipage}{.5\linewidth}
\begin{code}%
\>[0]\AgdaOperator{\AgdaFunction{\AgdaUnderscore{}+\AgdaUnderscore{}}}\AgdaSpace{}%
\AgdaSymbol{:}\AgdaSpace{}%
\AgdaDatatype{ℕ}\AgdaSpace{}%
\AgdaSymbol{→}\AgdaSpace{}%
\AgdaDatatype{ℕ}\AgdaSpace{}%
\AgdaSymbol{→}\AgdaSpace{}%
\AgdaDatatype{ℕ}\<%
\\
\>[0]\AgdaInductiveConstructor{zero}%
\>[7]\AgdaOperator{\AgdaFunction{+}}\AgdaSpace{}%
\AgdaBound{n}\AgdaSpace{}%
\AgdaSymbol{=}\AgdaSpace{}%
\AgdaBound{n}\<%
\\
\>[0]\AgdaInductiveConstructor{suc}\AgdaSpace{}%
\AgdaBound{m}%
\>[7]\AgdaOperator{\AgdaFunction{+}}\AgdaSpace{}%
\AgdaBound{n}\AgdaSpace{}%
\AgdaSymbol{=}\AgdaSpace{}%
\AgdaInductiveConstructor{suc}\AgdaSpace{}%
\AgdaSymbol{(}\AgdaBound{m}\AgdaSpace{}%
\AgdaOperator{\AgdaFunction{+}}\AgdaSpace{}%
\AgdaBound{n}\AgdaSymbol{)}\<%
\end{code}
\end{minipage}

\begin{code}%
\>[0]\AgdaFunction{+-comm}\AgdaSpace{}%
\AgdaSymbol{:}\AgdaSpace{}%
\AgdaSymbol{(}\AgdaBound{m}\AgdaSpace{}%
\AgdaBound{n}\AgdaSpace{}%
\AgdaSymbol{:}\AgdaSpace{}%
\AgdaDatatype{ℕ}\AgdaSymbol{)}\AgdaSpace{}%
\AgdaSymbol{→}\AgdaSpace{}%
\AgdaBound{m}\AgdaSpace{}%
\AgdaOperator{\AgdaFunction{+}}\AgdaSpace{}%
\AgdaBound{n}\AgdaSpace{}%
\AgdaOperator{\AgdaDatatype{≡}}\AgdaSpace{}%
\AgdaBound{n}\AgdaSpace{}%
\AgdaOperator{\AgdaFunction{+}}\AgdaSpace{}%
\AgdaBound{m}\<%
\\
\>[0]\AgdaFunction{+-comm}\AgdaSpace{}%
\AgdaInductiveConstructor{zero}%
\>[16]\AgdaInductiveConstructor{zero}%
\>[25]\AgdaSymbol{=}\AgdaSpace{}%
\AgdaInductiveConstructor{refl}\<%
\\
\>[0]\AgdaFunction{+-comm}\AgdaSpace{}%
\AgdaInductiveConstructor{zero}%
\>[16]\AgdaSymbol{(}\AgdaInductiveConstructor{suc}\AgdaSpace{}%
\AgdaBound{n}\AgdaSymbol{)}%
\>[25]\AgdaSymbol{=}\AgdaSpace{}%
\AgdaFunction{cong}\AgdaSpace{}%
\AgdaInductiveConstructor{suc}\AgdaSpace{}%
\AgdaSymbol{(}\AgdaFunction{+-comm}\AgdaSpace{}%
\AgdaInductiveConstructor{zero}\AgdaSpace{}%
\AgdaBound{n}\AgdaSymbol{)}\<%
\\
\>[0]\AgdaFunction{+-comm}\AgdaSpace{}%
\AgdaSymbol{(}\AgdaInductiveConstructor{suc}\AgdaSpace{}%
\AgdaBound{m}\AgdaSymbol{)}%
\>[16]\AgdaInductiveConstructor{zero}%
\>[25]\AgdaSymbol{=}\AgdaSpace{}%
\AgdaFunction{cong}\AgdaSpace{}%
\AgdaInductiveConstructor{suc}\AgdaSpace{}%
\AgdaSymbol{(}\AgdaFunction{+-comm}\AgdaSpace{}%
\AgdaBound{m}\AgdaSpace{}%
\AgdaInductiveConstructor{zero}\AgdaSymbol{)}\<%
\\
\>[0]\AgdaFunction{+-comm}\AgdaSpace{}%
\AgdaSymbol{(}\AgdaInductiveConstructor{suc}\AgdaSpace{}%
\AgdaBound{m}\AgdaSymbol{)}%
\>[16]\AgdaSymbol{(}\AgdaInductiveConstructor{suc}\AgdaSpace{}%
\AgdaBound{n}\AgdaSymbol{)}%
\>[25]\AgdaSymbol{=}\AgdaSpace{}%
\AgdaFunction{cong}\AgdaSpace{}%
\AgdaInductiveConstructor{suc}\AgdaSpace{}%
\AgdaSymbol{(}\AgdaFunction{trans}\AgdaSpace{}%
\AgdaSymbol{(}\AgdaFunction{+-suc}\AgdaSpace{}%
\AgdaBound{m}\AgdaSpace{}%
\AgdaBound{n}\AgdaSymbol{)}\AgdaSpace{}%
\AgdaSymbol{(}\AgdaFunction{+-comm}\AgdaSpace{}%
\AgdaSymbol{(}\AgdaInductiveConstructor{suc}\AgdaSpace{}%
\AgdaBound{m}\AgdaSymbol{)}\AgdaSpace{}%
\AgdaBound{n}\AgdaSymbol{))}\<%
\\
\>[0][@{}l@{\AgdaIndent{0}}]%
\>[2]\AgdaKeyword{where}%
\>[73I]\AgdaFunction{+-suc}\AgdaSpace{}%
\AgdaSymbol{:}\AgdaSpace{}%
\AgdaSymbol{∀}\AgdaSpace{}%
\AgdaBound{m}\AgdaSpace{}%
\AgdaBound{n}\AgdaSpace{}%
\AgdaSymbol{→}\AgdaSpace{}%
\AgdaBound{m}\AgdaSpace{}%
\AgdaOperator{\AgdaFunction{+}}\AgdaSpace{}%
\AgdaInductiveConstructor{suc}\AgdaSpace{}%
\AgdaBound{n}\AgdaSpace{}%
\AgdaOperator{\AgdaDatatype{≡}}\AgdaSpace{}%
\AgdaInductiveConstructor{suc}\AgdaSpace{}%
\AgdaSymbol{(}\AgdaBound{m}\AgdaSpace{}%
\AgdaOperator{\AgdaFunction{+}}\AgdaSpace{}%
\AgdaBound{n}\AgdaSymbol{)}\<%
\\
\>[.][@{}l@{}]\<[73I]%
\>[8]\AgdaFunction{+-suc}\AgdaSpace{}%
\AgdaInductiveConstructor{zero}%
\>[23]\AgdaBound{n}\AgdaSpace{}%
\AgdaSymbol{=}\AgdaSpace{}%
\AgdaInductiveConstructor{refl}\<%
\\
%
\>[8]\AgdaFunction{+-suc}\AgdaSpace{}%
\AgdaSymbol{(}\AgdaInductiveConstructor{suc}\AgdaSpace{}%
\AgdaBound{m}\AgdaSymbol{)}%
\>[23]\AgdaBound{n}\AgdaSpace{}%
\AgdaSymbol{=}\AgdaSpace{}%
\AgdaFunction{cong}\AgdaSpace{}%
\AgdaInductiveConstructor{suc}\AgdaSpace{}%
\AgdaSymbol{(}\AgdaFunction{+-suc}\AgdaSpace{}%
\AgdaBound{m}\AgdaSpace{}%
\AgdaBound{n}\AgdaSymbol{)}\<%
\end{code}
\begin{code}[hide]%
\>[0]\AgdaComment{--\ +-comm⁰\ :\ (m\ :\ ℕ)\ →\ m\ +\ zero\ ≡\ zero\ +\ m}\<%
\\
\>[0]\AgdaComment{--\ +-comm⁰\ zero\ =\ refl}\<%
\\
\>[0]\AgdaComment{--\ +-comm⁰\ (suc\ m)\ =\ cong\ suc\ (+-comm⁰\ m)}\<%
\\
%
\\[\AgdaEmptyExtraSkip]%
\>[0]\AgdaComment{--\ +-commˢ\ :\ (m\ n\ :\ ℕ)\ →\ m\ +\ suc\ n\ ≡\ suc\ (n\ +\ m)}\<%
\\
\>[0]\AgdaComment{--\ +-commˢ\ zero\ zero\ =\ refl}\<%
\\
\>[0]\AgdaComment{--\ +-commˢ\ zero\ (suc\ n)\ =\ cong\ suc\ (cong\ suc\ (sym\ (+-comm⁰\ \AgdaUnderscore{})))}\<%
\\
\>[0]\AgdaComment{--\ +-commˢ\ (suc\ m)\ zero\ =\ cong\ suc\ (+-commˢ\ m\ zero)}\<%
\\
\>[0]\AgdaComment{--\ +-commˢ\ (suc\ m)\ (suc\ n)\ =\ cong\ suc\ (trans\ (+-commˢ\ m\ (suc\ n))}\<%
\\
\>[0]\AgdaComment{--\ \ \ (trans\ (cong\ suc\ (sym\ (+-commˢ\ m\ n)))\ (+-commˢ\ (suc\ m)\ n)))}\<%
\\
%
\\[\AgdaEmptyExtraSkip]%
\>[0]\AgdaComment{--\ +-comm\ :\ (m\ n\ :\ ℕ)\ →\ m\ +\ n\ ≡\ n\ +\ m}\<%
\\
\>[0]\AgdaComment{--\ +-comm\ m\ zero\ =\ +-comm⁰\ m}\<%
\\
\>[0]\AgdaComment{--\ +-comm\ m\ (suc\ n)\ =\ +-commˢ\ m\ n}\<%
\end{code}

	\end{minipage}
	
	\vspace{-0.5\parskip}
	\hfill\begin{minipage}{0.95\textwidth}
	\light{\textbf{Remark 1.} Look at all the colors! \emoji{rainbow}}\\	
	\light{\textbf{Remark 2.} Is proving (\textit{m} + \textit{n} $\equiv$ \textit{n} + m) any different to (\textit{x} + \textit{y} $\equiv$ \textit{y} + \textit{x})? \emoji{thinking}}	
	\end{minipage}
	
	\textbf{Motivation}. If dependently-typed programs are proofs, and representing programs is essential to automating program synthesis, then \textbf{\textit{representing dependently-typed programs}} is key to \textbf{\textit{automating theorem proving}} (ATP).
		
	Two major issues in the literature:
	\begin{itemize}[topsep=-0.75\baselineskip]
		\item \textbf{Resource Uniformity.}\\
		Many ATP models/resources/interfaces for Coq, Lean.\\
		\textit{None for Agda}.
		\item \textbf{No Structural Fidelity.}\\
		Most ATP resources/frameworks today treat proofs as glorified text.\\
		\textit{Gone are all the colors. Names suddenly matter.}
	\end{itemize}

	{\sectionfont\vspace{\parskip}\textbf{Contributions (tl;dr)}}	
	\nfont
	
	\begin{itemize}[topsep=-0.25\baselineskip]
		\item \textbf{Machine Learning, for Agda}.\\
		We develop a package to faithfully extract the skeleton structure of dependently-typed program-proofs from type-checked Agda files. We apply the algorithm on Agda's public library ecosystem and release the result as a massive, highly elaborated ATP dataset.
		\item \textbf{Representation Learning for Dependent Types.}\\
		Capitalizing on this new resource, we present a representation learning model for expressions involving dependent types. Contra prior work, the model is structure-faithful, being invariant to $\alpha$-renaming, superficial syntactic sweetening, scope permutation, irrelevant definitions, \textit{etc}. 
	\end{itemize}
	
	\end{minipage}\hfill
	\begin{minipage}[t]{0.29\textwidth}
	\sectionfont{\textbf{Dataset}}	
	\nfont

	\textbf{Problem Generation}. Left-to-right language modeling assumes proving is a linear process. Truth begs to differ; the statement below is valid syntax:
	\begin{center}\begin{code}%
		\>[0]\AgdaFunction{+-comm}\AgdaSpace{}%
		\AgdaInductiveConstructor{zero}%
		\>[16]\AgdaSymbol{(}\AgdaInductiveConstructor{suc}\AgdaSpace{}%
		\AgdaBound{n}\AgdaSymbol{)}%
		\>[25]\AgdaSymbol{=}\AgdaSpace{}%
		\AgdaHole{\{!!\}}\<%
	\end{code}\end{center}

	\vspace{-0.5\parskip}
	\hfill\begin{minipage}{0.95\textwidth}
	\light{\textbf{Remark 3.} Proofs can have \textit{holes}: unfilled parts deferred for later. \emoji{cheese}}
	\end{minipage}
	
	We use Agda's type-checker to find \textit{all possible holes} in all written proofs. For each hole, we record the \textit{goal type} (the type of the hole) and the \textit{typing context} (all proven premises currently available). Ground truth corresponds to a selection (and arrangement) of the context (how to fill the hole).
	
	\vspace{-0.5\parskip}
	\hfill\begin{minipage}{0.95\textwidth}
	\light{\textbf{Remark 4.} Correct \textbf{premise selection} goes a long way towards ATP. \emoji{jigsaw}}
	\end{minipage}
	
	We export the extracted problems not \textit{only} as \textbf{strings}, but \textit{also} as \textbf{structures}.\linebreak
	The export preserves and specifies all type information available to the checker, including \textbf{references} and \textbf{token structure at the subtype level}.
	
	
	Post-tokenization, this is what the \textit{types} of \AgdaDatatype{$\mathbb{N}$}, \AgdaFunction{$+$} and \AgdaFunction{+-comm} look like:
	
	\begin{tikzpicture}[
			ast/.style={draw,gray,thick,outer sep=2pt,opacity=0.5},
			lab/.style={font=\large},
			ref/.style={->,dashed,gray,opacity=0.4,thick},
			var/.style={gray,opacity=0.8,execute at begin node={\strut(},execute at end node=)},
			t/.style={text height=1.5ex, text depth=.25ex, rectangle, outer sep=0pt},
			scale=1.45, transform shape]
		\tikzset{level 1/.style={level distance=22pt}, edge from parent/.append style={thick}}
		\tikzfont
		\begin{scope}[xshift=1cm]
		\Tree
		[.\node(t11) {•};
				\node (t12) {\AgdaPrimitive{Set}};
		]
		\node[ast,fit=(t11)(t12)](nat) {};
		\node[lab,xshift=5.5pt,yshift=4pt]() at(nat.north west) {\AgdaDatatype{ℕ}};
		\begin{scope}[xshift=3.5cm,yshift=-1.5cm]
		\Tree
		[.\node (t21) {•};
			[.{\sanode{}}
				\node[var] (plus-l) {\AgdaDatatype{ℕ}};
				[.{\sanode{}}
					\node[var] (plus-r) {\AgdaDatatype{ℕ}};
					\node[var] (plus-out) {\AgdaDatatype{ℕ}};
				]
			]
		]
		\node[ast,fit=(t21)(plus-l)(plus-out)] (+) {};
		\node[lab,xshift=4.5pt,yshift=4pt]() at(+.north west) {\AgdaFunction{+}};
		\begin{scope}[xshift=6cm,yshift=3.2cm]\Tree
		[.\node (t31) {•};
			[.{\dpnode{}}
				\node[var] (m-def) {$\AgdaBound{m}:\AgdaPrimitive{Set}$};
				[.{\dpnode{}}
					\node[var] (n-def) {$\AgdaBound{n}:\AgdaPrimitive{Set}$};
					[.{\fanode{}}
						[.{\fanode{}}
							{(\AgdaPrimitive{≡})}
							[.{\fanode{}}
								[.{\fanode{}}
									\node[var] (+-ref-1) {\AgdaFunction{+}};
									\node[var] (m-ref-1) {\AgdaBound{m}};
								]
								\node[var] (n-ref-1) {\AgdaBound{n}};
							]
						]
						[.{\fanode{}}
							[.{\fanode{}}
								\node[var] (+-ref-2) {\AgdaFunction{+}};
								\node[var] (n-ref-2) {\AgdaBound{n}};
							]
							\node[var] (m-ref-2) {\AgdaBound{m}};
						]
					]
				]
			]
		]
		\node[ast,fit=(t31)(m-def)(+-ref-1)(m-ref-2)] (comm) {};
		\node[lab,xshift=19pt,yshift=4pt]() at(comm.north west) {\AgdaFunction{+-comm}};
		\end{scope}\end{scope}\end{scope}
		\draw[ref] (plus-l) -- (nat);
		\draw[ref] (plus-r) -- (nat);
		\draw[ref] (plus-out) -- (nat);
		\draw[ref] (m-def) -- (nat);
		\draw[ref] (n-def) -- (nat);
		\draw[ref] (+-ref-1) -- (+);
		\draw[ref] (+-ref-2) -- (+);
		\draw[ref] (m-ref-1) -- (m-def);
		\draw[ref] (m-ref-2) -- (m-def);
		\draw[ref] (n-ref-1) -- (n-def);
		\draw[ref] (n-ref-2) -- (n-def);
	\end{tikzpicture}

%		\vspace{-0.5\parskip}
		\hfill\begin{minipage}{0.95\textwidth}\light{\textbf{Remark 5.} Note the AST and referencing structures \emoji{christmas-tree}\emoji{christmas-tree}\emoji{christmas-tree}}\\
		\light{\textbf{Remark 6.} Contrast with the tokenization of GPT-4o below. \raisebox{-0.25em}{\rotatebox{45}{\emoji{straight-ruler}}}}
		\end{minipage}
		
		\newcommand{\cbp}[1]{\framebox{\strut#1}}
		\newcommand{\cbg}[1]{\framebox{\strut#1}}
		\newcommand{\cbr}[1]{\framebox{\strut#1}}
		\newcommand{\cbb}[1]{\framebox{\strut#1}}
		\newcommand{\cby}[1]{\framebox{\strut#1}}
		{\tikzfont\centering
		\cbp{data}\cbg{~}\cby{$\mathbb{N}$}\cbr{ :}\cbb{ Set}\cbp{<newline>}\cbg{\_}\cby{+}\cbr{\_}\cbb{ :}\cbg{$\mathbb{N}$}\cby{$\to$}\cbr{}\cbb{$\mathbb{N}$}\cbr{<newline>}\cbb{+-}\cbp{comm}\cbg{ :}\cby{ (}\cbr{m}\cbb{ n}\cbp{ :}\cbg{}\cby{$\mathbb{N}$}\cbr{)}\cbb{ $\to$}\cbp{ m}\cbg{ +}\cby{ n}\cbr{}\cbb{$\equiv$}\cbp{ m}\cbg{ +}\cby{ m}\cbr{<newline>}
		}
	\end{minipage}\hfill
	\begin{minipage}[t]{0.29\textwidth}
	\sectionfont{\textbf{Representation Learning}}
	\nfont

	We build representations for lemmas and holes on the basis of their types.
	
	\textbf{Architecture.} We use a fully-attentive bidirectional Transformer encoder, where full attention is restricted to tokens within the same type, augmented with various representational adjustments.
	\begin{enumerate}[topsep=-0.75\baselineskip]
		\item \textbf{Tree PE}. We use positional encodings that employ an inductive parameterization of the group structure of binary branching trees. These relieve the model from having to ``parse'' the type's symbolic sequentialization.
		\item \textbf{Variable Binding.} We resolve nominal indexing, and represent variable references by the representation of the reference's path relative to the binder.
		\item \textbf{Scope Referencing.} We organize lemmas into a \textsc{poset} according to their dependency levels. We then build representations in dependency-sorted minibatches, and represent lemma references by the representations of their referents.
		\light{\hfill (here: $\mathbb{N}$ < $+$ < +-comm)}
		\item \textbf{Efficient Attention.} We use linear attention combined with a Taylor-approximation of the exponential map to efficiently avoid the quadratic explosion -- without losing expressivity.
	\end{enumerate}
	
	\vspace{\baselineskip}
	\textbf{Training.} We train with infoNCE in a premise selection setup using a subset of Agda's standard library, and evaluate in proximal and distant domains.
	
	{\newcommand{\stdlib}{\texttt{stdlib}}
	\newcommand\unimath{\texttt{Unimath}}
	\newcommand\typetopo{\texttt{TypeTopo}}
	\newcommand\agdaQuill{\textsc{Quill}}	

	\newcommand{\sstat}[2]{\hphantom{1}\ensuremath{#1}}
	\newcommand{\bstat}[2]{\ensuremath{#1}}
	\newcommand{\bfstat}[2]{\ensuremath{\mathbf{#1}}}
	
	\newcolumntype{C}{>{\centering\arraybackslash}X}
	\newcolumntype{L}{>{\raggedright\arraybackslash}X}
	\newcolumntype{R}{>{\raggedleft\arraybackslash}X}
	
	\centering
	\begin{tabularx}{0.95\textwidth}{
		@{}l@{\quad}L@{~/~}R@{\quad}R@{~/~}L@{\quad}R@{~/~}L@{\quad}R@{~/~}L@{}
%		@{}l@{~~}c@{~~}c@{\hskip 36.5pt}c@{~~}c@{\hskip 36.5pt}c@{~~}c@{\hskip 36.5pt}c@{~~}c
	}
	& \multicolumn{8}{c}{Average and R-Precision}\\
	\multicolumn{1}{@{}l}{\textsc{Model}}
	& \multicolumn{2}{@{\quad}C@{\quad}}{\stdlib{}:\textsc{id}}
	& \multicolumn{2}{C@{\quad}}{\stdlib{}:\textsc{ood}}
	& \multicolumn{2}{@{\quad}C@{\quad}}{\unimath}
	& \multicolumn{2}{@{\quad}C@{\quad}}{\typetopo}\\
	\toprule
	~\agdaQuill 	&
		\bfstat{50.2}{0.5} & \bfstat{40.3}{0.4} & \bfstat{38.7}{0.7} & \bfstat{31.1}{0.9} & \bfstat{27.0}{0.4} & \bfstat{17.4}{0.3} & \bfstat{22.5}{0.3} & \bfstat{15.4}{0.5} \\
	~~- (4)		& \bstat{47.0}{0.4} & \bstat{36.2}{0.5} & \bstat{37.1}{0.5} & \bstat{29.2}{0.6} & \bstat{26.8}{0.2} & \bstat{17.0}{0.1} & \bstat{21.4}{0.5} & \bstat{14.4}{0.4} \\
	~~- (1)	 & \bstat{44.5}{1.5} & \bstat{34.1}{1.8} & \bstat{30.7}{0.4} & \bstat{24.0}{0.4} & \bstat{24.8}{3.4} & \bstat{15.5}{2.9} & \bstat{18.8}{1.1} & \bstat{12.3}{0.9}\\
	~~- (2)	& \bstat{35.8}{2.7} & \bstat{25.9}{2.7} & \bstat{25.5}{2.7} & \bstat{19.1}{2.4} & \bstat{19.7}{1.8} & \bstat{11.6}{1.5} & \bstat{17.7}{3.0} & \bstat{11.0}{3.1}\\	
	~Transformer & \bstat{10.9}{0.4} & \sstat{3.7}{0.2} & \sstat{8.5}{0.2	} & \sstat{4.5}{0.1} & \sstat{9.4}{0.3} & \sstat{3.9}{0.1} & \sstat{5.8}{0.0} & \sstat{0.9}{0.0} \\
	\bottomrule	
	\end{tabularx}
	}
	
	\hfill\begin{minipage}{0.95\textwidth}
		\light{\textbf{Remark 7.} Structural adjustments >> architectural adjustments. \emoji{balance-scale}}
	\end{minipage}


	\sectionfont
	\vspace{\parskip}
	\textbf{Learn more}
	\nfont
	
	For more details, take a look at:
	\begin{itemize}[topsep=-0.75\baselineskip]
		\item \url{agda.readthedocs.io}\\
		for an intro to Agda
		\item \url{github.com/omelkonian/agda2train}\\ for the proof extraction code 
		\item \url{github.com/konstantinosKokos/quill}\\ for the Python interface and neural engine
		\item \url{arxiv.org/abs/2402.02104}\\ for prose, figures, tables with numbers, etc.
	\end{itemize}

	\vspace{\parskip}\vspace{\parskip}\vspace{\parskip}\vspace{\parskip}
		
	\includegraphics[height=0.2\textwidth]{cost.pdf}\hspace{1.5em}
	\raisebox{1.25em}{\includegraphics[height=0.12\textwidth]{eu.png}}
	
	\end{minipage}%

\end{document}

